\section{Tax information before tax calculation}
\label{sec:shared-tax-information}

The mixed book has recorded profit, but it still lacks a tax assessment.
The earlier return-only reference does contain a tax-close method. Whether
that method can be reused is a substantive question: the new book contains
lease rights and a modification gain that the old book did not represent.
The unchanged method refuses the mixed state because its closed state shape
does not admit those additional accounts and lease records. The two original
tax examples still reproduce their complete results. Thus the refusal identifies
an incompatible extension, rather than a regression in the original method.

\subsection{A new income category changes the tax starting point}

The original pretax function sums the original income accounts, excluding
current and deferred tax. If that function is applied diagnostically to the
mixed balances, it returns a loss of 15. The complete recorded report instead
includes the lease-modification gain of 20. Under the signed-account convention,
the additional \texttt{other\_gains} balance is $-20$, so
\[
  P_{\mathrm{old\ accounts}}=-15,
  \qquad
  P_{\mathrm{mixed}}=-15-(-20)=5.
\]
The difference is an omitted income category, not a rounding error. Replacing
the state validator while leaving the original income list unchanged would
therefore calculate the wrong pretax starting point for the stated mixed-book
target. It would not preserve the meaning of the old calculation.

Even the corrected starting point does not determine tax. The return example
has a signed reconciliation adjustment of $-24$. Mechanical reuse of that
adjustment and the original illustrative rate would give
\[
  \frac{-15-24}{4}=-\frac{39}{4},
  \qquad
  \frac{5-24}{4}=-\frac{19}{4}.
\]
These expressions diagnose the effect of the omitted gain; neither is an
assessed tax or a cash refund. The lease deduction timing, the tax treatment
of the modification, the remaining right-of-use tax bases and the treatment
of a loss have not been supplied. The original finite method refuses its
negative-current-tax route. Copying the rate or reversing every deferred-tax
balance would not supply the missing facts.

Dates create a separate incompatibility. The original close on 31 December
stores 1 January as an exclusive tax endpoint. The mixed-book invariant
currently requires that endpoint to be no later than the state date. Those
two meanings cannot be combined by changing only the endpoint field. Moreover,
the mixed example's execution window runs from 1 January 2024 to 1 January 2026:
731 days. That window contains the selected lease and return events, but does
not establish a single annual tax interval or complete fiscal activity.

\subsection{Keep missing declarations and rejected premises visible}

A separate information assessor now accepts the unchanged recorded-report
request, an explicitly stated tax interval and identified information
declarations. The interval distinguishes its start, its last included calendar
day and its exclusive endpoint. The first finite implementation requires the
start to equal the recorded opening and every requested event to lie inside
the interval. The last recorded event may precede the asserted coverage date;
that assertion remains a declaration, not evidence that all activity is known.

For this mixed example, the assessor identifies 61 required declarations. They
cover jurisdiction and regime, the interval basis and coverage findings, rate,
loss and cash-settlement treatment, deferred recognition and recovery, every
pretax account's treatment and signed reconciliation adjustment, balance-sheet
tax bases, and the tax bases and timing of each lease right. Zero balances
remain in the account population. Both lease rights A and B remain relevant:
the native state retains fully terminated B with zero asset and liability,
and its earlier modification still requires a stated treatment.

Each supplied value carries an assumption or source-reference description.
Each unknown carries its reason. A source reference is retained as supplied;
the assessor does not authenticate it. With no information supplied, all
61 declarations remain unresolved. Removing any one from the hypothetical
complete control again leaves an unresolved requirement. A false coverage or
recognition finding remains false and is reported as unsupported. Extra fields,
wrong scalar types, forged request identities and incompatible interval
boundaries are invalid.

The complete native report accompanies the assessment. Its stopped event
prefix, component memory, gross cash movements and unresolved annual coverage
are preserved. When a requested payment fails, supplying every information
declaration does not erase that failure: completion of the requested sequence
remains unresolved. A first-year interval can be stated separately when its
requested events fit that interval, without relabelling the two-year example.

\subsection{Supplied information is not an executable tax rule}

The positive result is deliberately called \texttt{information\_supplied}.
It means the finite declaration set is present and has the required types.
The hypothetical controls exercise that representation; their supplied zero
bases and adjustments are not inferred from the financial statements. A text
field describing a deduction can still fail to specify a valid calculation.
Likewise, individually typed amounts can be inconsistent when aggregated.

Accordingly, every result keeps calculation readiness, source and accounting
admission, tax execution and annual-close readiness false. No tax amount is
posted. The next calculation design must state how account adjustments,
right-level tax bases, deferred recognition and settlements relate, and must
check their joint consistency. This preserves useful progress: missing facts
and competing declarations can be represented and examined now, while the
remaining numerical and accounting questions remain explicit.

\subsection{Reconcile a declared current-tax calculation}

The next finite calculation makes one further assumption explicit: a single
nonnegative rate applies to nonnegative taxable profit before credits or any
other regime-specific adjustments. This is a hypothetical arithmetic model,
named \texttt{uniform\_rate\_no\_credits\_v1}. It does not select a country's
tax rules. The author must select that model affirmatively; a false selection
is unsupported and an unknown selection remains unresolved.

Each declared account adjustment is now supported by identified allocation
items. An item taken from the completed return history must retain its original
event identity and signed amount, and must appear exactly once. Additional
items have explicit assumption references. The items assigned to an income
account must sum to that account's declared adjustment. Their category assignment
is supplied by the author; the arithmetic check does not establish its legal
appropriateness. In particular, the native return adjustment is not added again
after it has already been included in those account totals.

There are also four cross-field basis checks. At both opening and closing,
the declared asset bases of all lease rights must sum to the aggregate
right-of-use tax basis, and the declared liability bases must sum to the
aggregate lease-liability basis. The zero-valued right B participates in those
checks. A mismatch rejects the conditional calculation while preserving the
complete recorded financial report and the inconsistent declarations for
inspection. Neither carrying amount nor termination status supplies a tax basis.

\begin{samepage}
Let $P$ be recorded pretax income over all seven income accounts, $a_i$ the
identified signed adjustment items and $r$ the explicitly supplied rate. The
conditional taxable-profit expression and gross current tax are
\[
 T=P+\sum_i a_i,
 \qquad
 C=rT\quad\text{when }T\geq0.
\]
\end{samepage}
Here $C$ is gross conditional current tax under the selected arithmetic rule.
It is not an incremental expense, accrued payable or cash payment. The
calculation does not subtract opening accruals, apply credits, recognize
deferred tax or change the recorded closing state.

For the mixed example, $P=5$ and the retained return adjustment is $-24$.
Without any additional declared item, $T=5-24=-19$. The result is
\texttt{loss\_route\_required}; no negative tax amount or refund is computed.
A text declaration requesting a refund cannot change that result because it
does not supply an executable loss rule. Two hypothetical controls illustrate
the nonnegative branch. An expressly supplied additional adjustment of 19 gives
$T=0$ and $C=0$. An additional adjustment of 40 gives $T=21$ and, at the
illustrative rate $1/4$, $C=21/4$. Those extra adjustments are test assumptions,
not inferred company facts or recommended tax treatments.

The comparison retains the entire financial report, including cash of 361 and
recorded profit of 5, in every accepted arithmetic result. Inconsistent allocation
or basis declarations retain that same report alongside their rejection. Tests
also change native adjustment amounts, duplicate or omit identities, alter the
model finding and tamper with computed results. These checks establish the
specified arithmetic and preservation properties. A loss-utilization rule,
tax settlement, deferred recognition and annual close remain separate work.

\subsection{Carry a loss quantity without inventing a tax asset}

The negative expression of $-19$ needs a rule before it can affect a later
tax calculation. The finite extension supplies a conditional loss ledger.
Each lot has an identity, an origin reference, an availability date, an explicit
expiry declaration, original and remaining quantities, cumulative usage and
expiry, and an eligibility finding. An unknown finding remains unresolved;
a false finding is unsupported. A generated loss therefore becomes a quantity
in this separate ledger only under an affirmative, traceable assumption.
It is not a cash refund or a recognized deferred-tax asset.

The selected illustrative policy evaluates all declared usage on the last
included calendar day of an interval. An exclusive expiry on that day makes
the lot unavailable for usage. Expiry at the following endpoint permits usage
on the preceding day, but the remainder expires before the next opening.
No expiry is itself an explicit assumption. These conventions specify this
finite model; another jurisdiction or an intraperiod method may require a
different rule. Neither FIFO nor a legal utilization percentage is assumed.
The author supplies the chosen lot allocations and an interval usage cap.

For each lot, original loss equals remaining, used and expired quantities.
Across a step, with $L_0$ and $L_1$ denoting opening and closing remaining
quantities, $G$ generation, $U$ usage and $E$ expiry,
\[
 L_0+G=U+E+L_1.
\]
All quantities on this line are nonnegative exact rationals. A generated loss
becomes available at the exclusive endpoint and cannot be used in its own
generation interval. A usage cannot exceed its lot's remaining quantity,
the declared interval cap or positive taxable profit. Zero and exhausted lots
remain identified, preventing the same identity from being introduced again.
Consecutive ledger entries must agree on both interval endpoints and their
opening and closing quantities.

The checked negative example generates 19 under explicit eligibility and
no-expiry premises. A separate hypothetical 2026 calculation supplies taxable
profit of 21. Explicit use of 19 leaves taxable profit of 2 and gross conditional
tax of $1/2$ at rate $1/4$. This connects the loss ledger across periods; it does
not establish financial-book carry. The 2026 financial opening is independently
declared for the test and is expressly not the carried closing state of the
2024--2026 example. Both complete conditional calculations remain attached.

Partial use makes the expiry rule consequential. Using 7 of 19 leaves 12 when
the expiry is later than the next endpoint. If the exclusive expiry equals that
endpoint, 7 may be used on the preceding day and the remaining 12 expires.
If expiry instead equals the usage day, the usage is rejected; with no usage,
all 19 expires. These outcomes follow the selected policy, rather than a
silent decision to use the oldest lot or to preserve an expired balance.

The operation proposes a complete next loss state and per-lot movements.
Applying it recomputes the proposal and compares the complete result, including
the original financial calculation. Changed amounts, origins, eligibility,
scope, cursors or even Boolean-to-integer substitutions are rejected. This is
a pure in-memory application: it creates no durable commit, tax journal,
deferred asset or annual authorization. Connecting the conditional tax amount
and loss ledger to an actual financial tax balance still requires separately
specified recognition, accrual and settlement methods.

\subsection{Accrue current tax before choosing its payment}

A conditional calculation supplies an amount, but recognition and payment
change different financial balances. The next finite profile makes those
changes explicit. It begins with the complete recorded return-and-lease
calculation, recomputes the conditional loss result and requires an affirmative
recognition premise. An unknown premise remains unresolved; a false premise
does not produce a financial state. The profile applies only when the native
period has no current-tax expense already recorded and its tax-close cursor
still equals the assessment start. It therefore does not guess how an opening
tax accrual or an earlier assessment should be allocated.

The successor adds a separate income-tax payable to the original 25 accounts.
That liability is not trade payable. If the recomputed gross conditional tax
after loss usage is $C$, the accrual debits current-tax expense and credits
income-tax payable by $C$. Cash does not change. A zero amount produces a zero
journal. Existing deferred-tax assets, deferred-tax liabilities and deferred-tax
expense remain exactly as recorded; the loss balance does not create a new
deferred asset.

Payment requires its own identity, date, positive amount and affirmative
instruction. It debits income-tax payable and credits cash, subject to both the
outstanding liability and available cash. The author supplies the amount;
neither full payment nor a refund is inferred. A shortage of cash blocks that
payment while leaving a valid accrual available. This illustrative profile
classifies its specified tax payments as operating cash movements.

\begin{samepage}
Let $S$ be cumulative payment against the newly accrued amount $C$, with
$0\leq S\leq C$. Relative to the original recorded statements, the changes are
\[
 \Delta\text{profit}=-C,\qquad
 \Delta\text{tax payable}=C-S,\qquad
 \Delta\text{cash}=-S.
\]
\end{samepage}
Assets decrease by $S$, liabilities increase by $C-S$ and equity decreases by
$C$, preserving the balance-sheet identity. In the indirect operating cash
bridge, the lower profit and the increase in tax payable contribute
$-C+(C-S)=-S$. This equals the separately recorded cash payment. Paying the
liability does not charge the same expense again.

For the separate hypothetical 2026 calculation, loss usage leaves conditional
tax of $1/2$. Accrual records expense and payable of $1/2$ while cash stays at
500. An explicit payment of $1/4$ leaves cash of $1999/4$ and payable of $1/4$;
a second payment of $1/4$ clears the payable and leaves cash of $999/2$.
The recorded deferred-tax asset of 10 and liability of 4 are unchanged.
The original mixed example instead generates the loss quantity 19 and a zero
tax journal. These are consequences of the declared synthetic inputs, not
assessments of either a company or a jurisdiction.

The financial date and assessment coverage have separate meanings. Accrual is
dated on the last included day; coverage ends at the following exclusive
endpoint. The successor records that endpoint separately and leaves the native
tax-close cursor unchanged. Payments cannot precede the last tax operation or
fall outside the original financial period. An end-of-period accrual therefore
cannot silently accept a payment in the next year: financial-period carry must
first be specified.

The original native schema refuses the additional account and the exclusive
close cursor, so this is a distinct financial profile. Its validation replays
the retained calculation and every tax payment, then compares the complete
state. Changed balances, dates, loss results, payment histories and proposal
flags are rejected. Repeating the same accrual request produces the same pure
proposal; it does not establish globally unique durable posting. Later ordinary
return or lease events, financial-year carry, actual accounting admission and
annual completeness remain separate requirements.
